\chapter{Methodology and Theory}

This chapter outlines the theoretical foundations of explicit dynamics and the specific mechanical principles governing the EN 397 drop test simulation.

\section{Explicit Dynamics Theory}
Explicit dynamics apply the central difference method to solve kinematic equations of motion directly without global stiffness matrix inversions. This is highly efficient for high-nonlinearity, millisecond-duration impact events.

\subsection{Central Difference Integration}
The semi-discrete equations of motion at time $t_n$ are expressed as:
\begin{equation}
    M \ddot{u}_n = F^{ext}_n - F^{int}_n
\end{equation}
where $M$ is the diagonal (lumped) mass matrix, $\ddot{u}_n$ is the nodal acceleration vector, $F^{ext}_n$ represents external applied loads (including contact forces), and $F^{int}_n$ represents internal nodal forces.

LS-DYNA employs a modified central difference time integration scheme (Leapfrog method) to advance the solution~\cite{lsdyna2021}. The nodal velocities and displacements are updated according to:
\begin{equation}
    \dot{u}_{n+1/2} = \dot{u}_{n-1/2} + \ddot{u}_n \Delta t_n
\end{equation}
\begin{equation}
    u_{n+1} = u_n + \dot{u}_{n+1/2} \Delta t_{n+1/2}
\end{equation}
where $\Delta t_{n+1/2} = \frac{1}{2}(\Delta t_n + \Delta t_{n+1})$. This explicit approach allows for efficient processing of complex contact interactions between the rigid headform, the EPS liner, the polymer shell, and the rigid ceiling.

\subsection{Stability and Time Step Control}
The central difference method is conditionally stable. The time step $\Delta t$ is strictly limited by the Courant-Friedrichs-Lewy (CFL) condition~\cite{belytschko2013}:
\begin{equation}
    \Delta t \leq \Delta t_{crit} = \frac{L_{min}}{c}
\end{equation}
where $L_{min}$ is the characteristic length of the smallest element in the mesh and $c$ is the local speed of sound in the material ($c = \sqrt{E/\rho}$). For this study, the 2 mm mesh resolution at the crown impact zone dictates the global time step, ensuring that the stress wave propagation is accurately captured during the 49 J impact.

\section{Unit System and Numerical Consistency}
A consistent [kg, mm, s] unit set was adopted for LS-DYNA stability. Strength and contact pressure are reported in kg/mm$^2$ ($1 \text{ kg/mm}^2 \approx 9.81 \text{ MPa}$), while Young’s ($E$) and Tangent ($E_{tan}$) Moduli are in GPa. Transmitted forces and impact energy use kN and J, respectively.
To optimise the computational domain and stabilise the explicit penalty contact 
algorithms, a kinematic shift was implemented based on the principle of 
Galilean Invariance. Rather than a transient drop of a 5.0 kg mass onto a 
stationary helmet, the rigid ``floor'' was modelled as a stationary ceiling. 
The helmet and the 5.0 kg spherical headform assembly were grouped and 
assigned a uniform upward initial velocity of 4.43 m/s. Figure~\ref{fig:assembly_setup} 
illustrates the numerical assembly and the kinematic boundary conditions.

\begin{figure}[H]
    \centering
    \begin{subfigure}[b]{0.45\textwidth}
        \includegraphics[width=\textwidth]{cad_impact_setup.png}
        \caption{EN 397 Top Impact}
        \label{fig:top_setup}
    \end{subfigure}
    \hfill
    \begin{subfigure}[b]{0.45\textwidth}
        \includegraphics[width=\textwidth]{model_30_oblique_setup.png}
        \caption{30$^\circ$ Oblique Impact}
        \label{fig:oblique_setup}
    \end{subfigure}
    \caption{Numerical assembly and impact configurations. The Galilean upward velocity shift is applied to the headform-helmet assembly against a stationary rigid boundary.}
    \label{fig:assembly_setup}
\end{figure}

According to classical mechanics, the laws of motion are invariant in all 
inertial reference frames related by a constant velocity shift~\cite{goldstein2002}. 
In explicit finite element analysis, this transformation is a standard 
technique to isolate the impact event from free-flight processing time, 
provided that the element formulations are objectively invariant 
\cite{lsdyna2021, belytschko2013}. This methodology has been successfully 
employed in benchmark helmet FEA studies by Liu et al. \cite{liu2003} and 
more recently in moving-load impact scenarios by Høgsberg and Pedersen 
\cite{hogsberg2017}. 

The physics of the 49 Joules impact ($E_k = \frac{1}{2}mv^2$) remains 
mathematically identical at the moment of contact. Because the 
impact duration in these simulations is extremely short ($\approx 8$--10 ms), 
the additional acceleration due to gravity ($g = 9.81$ m/s$^2$) during the 
contact phase contributes less than 0.2\% to the total impulse and can be 
safely neglected in favor of contact stability and domain optimization.

\section{Contact Mechanics and the Penalty Formulation}
The interaction between the helmet shell, the EPS liner, and the rigid headform was governed by LS-DYNA's \texttt{*AUTOMATIC\_SURFACE\_TO\_SURFACE} algorithm. To prevent non-physical nodal penetration, this algorithm employs a penalty-based formulation. When a slave node penetrates a master segment, an elastic compression-only spring is dynamically instantiated between them to resist the penetration~\cite{belytschko1992}.

The contact force vector, $\mathbf{F}_c$, applied to the penetrating node is proportional to the penetration depth, $\delta$:
\begin{equation}
    \mathbf{F}_c = k_c \cdot \delta \cdot \mathbf{n}
\end{equation}
where $\mathbf{n}$ is the unit normal vector of the master segment, and $k_c$ is the interface contact stiffness. For solid elements, LS-DYNA derives this stiffness based on the bulk modulus of the interacting materials:
\begin{equation}
    k_c = \frac{f_s \cdot K \cdot A}{V}
\end{equation}
where $K$ is the bulk modulus of the slave element, $A$ is the segment face area, $V$ is the element volume, and $f_s$ is the penalty scale factor~\cite{lsdyna_theory}. 

In this investigation, a penalty scale factor ($f_s$) of 0.1 was utilized. This value optimizes the balance between enforcing kinematic constraints (preventing unphysical mesh overlap) and maintaining thermodynamic stability. Higher penalty factors artificially stiffen the contact interface, potentially driving the required explicit time-step below the Courant-Friedrichs-Lewy (CFL) stability limit and inducing catastrophic high-frequency noise.
